\section{Rooted labeling interfaces and path inputs}
\label{common:foundations}

All graphs in this manuscript are finite. A path length is its number
of edges, while $P_n$ denotes a path on $n$ vertices. These conventions
are especially important when a path $P_L$ is joined to a hub by a new
edge: the new arm then has length $L$.

\begin{definition}
Let $T$ be a tree with $m$ edges. A \emph{graceful labeling} is a
bijection $f:V(T)\to\interval{0}{m}$ such that the edge differences
$\abs{f(u)-f(v)}$ are precisely $1,\ldots,m$, each once. It is an
\emph{$\alpha$-labeling of integer index $K$} if every edge has one endpoint
of label at most $K$ and the other of label greater than $K$,
where $0\le K\le m$. In particular, the singleton has index zero.
For a distinguished root $r$, \emph{root-zero} means $f(r)=0$.
\end{definition}

The rooted spider $\Sp{L_1,\ldots,L_d}$ is obtained by identifying
one endpoint of each of $d$ otherwise disjoint paths of positive edge
lengths $L_i$. Its distinguished root is the common hub. We also allow
$d\le2$, and allow $d=0$ for the one-vertex tree. Repeated lengths are
written $L^t$ when the multiplicity $t$ is specified.

\begin{definition}
For a rooted $m$-edge tree $T$, write $k\in\Om(T)$ if there is a
labeling with root zero, nonroot labels exactly
$\interval{k+1}{k+m}$, and differences exactly $\interval{1}{m}$.
Only integers $k\ge0$ are considered. A realization at $k=0$ is a
hub-zero graceful labeling when the root is a spider hub.
\end{definition}

\begin{lemma}[Forced index]\label{common:root-index}
A root-zero $\alpha$-labeling of a spider has index
\[
 K=\sum_{i=1}^d\floor{L_i/2}.
\]
\end{lemma}
\begin{proof}
The root belongs to the lower class. Its bipartition class consists of
the root and the vertices at positive even distance along the arms,
and hence has $1+\sum_i\floor{L_i/2}$ vertices. All labels from zero
through the index occur in this class, and every label occurs once.
Its size is therefore $K+1$.
\end{proof}

The classical graceful-tree problem asks for a graceful labeling of every tree. Spiders provide a natural family with one branching vertex. The published five-leg theorem is due to Panpa, Imnang and Wasuanankul~\cite{Panpa2025}; the six-leg rooted-offset approach appears in the 2026 preprint of Rhys~\cite{Rhys2026}. Related constructions appear in Bahls, Lake and Wertheim~\cite{BahlsLakeWertheim2010}, Shan and Zhong~\cite{ShanZhong2026}, and Dumitru and Nacu~\cite{DumitruNacu2026}. The present proofs use the classical path and amalgamation results below, and give the subsequent construction arguments explicitly. The six-leg conclusion itself is not a premise.\par\medskip

\subsection{Classical path results with precise normalizations}

We use the following established path theorem. The $\alpha$-statements
are due to Rosa and Cattell \cite{Rosa1977,Cattell2007}; the three
forms are restated as Theorems~2.1--2.3 of \cite{Panpa2025}.

\begin{theorem}[Prescribed labels on paths]\label{common:path-classical}
For $n\ge2$ the following hold.
\begin{enumerate}
\item Any specified vertex of $P_n$ can have label zero in an
$\alpha$-labeling, except the central vertex of $P_5$.
\item An endpoint of $P_n$ can have any specified graceful label in
$\interval{0}{n-1}$.
\item An endpoint can have any such label in an $\alpha$-labeling,
except when $n=4s+1$, $s\ge1$, and the label is $s$ or $3s$.
\end{enumerate}
The one-vertex path is treated separately with its unique label zero.
\end{theorem}

In construction formulas we often normalize a path to labels
$\interval{1}{n}$ instead. This adds one to every label and to the
cut, preserves differences, and changes the two exceptional labels
to $s+1$ and $3s+1$.

\begin{corollary}[Rooted pairs]\label{common:alpha-pair}
A rooted pair $\Sp{a,b}$ has a root-zero $\alpha$-labeling unless
$(a,b)=(2,2)$. Its index, when it exists, is
$\floor{a/2}+\floor{b/2}$. Every single arm is root-zero $\alpha$.
\end{corollary}
\begin{proof}
The pair is $P_{a+b+1}$ with the distinguished vertex at distances
$a,b$ from its endpoints. It is the exceptional rooted $P_5$ exactly
when $a=b=2$. A rooted single arm is covered by the same theorem, or
by the Walecki word $0,L,1,L-1,2,L-2,\ldots$.
\end{proof}

Every rooted pair is nevertheless ordinary hub-zero graceful. The only
case not supplied by the alpha statement is $\Sp{2,2}$, for which the
two outward words $(4,1)$ and $(2,3)$ give the colors $4,3,2,1$.

\begin{lemma}[Even-order path interface]\label{common:even-path}
For $r\ge1$ and $1\le v\le r$, there is an $\alpha$-labeled path
on $\interval{1}{2r}$, cut at $r$, starting at $r+v$ and ending at
$v$. There is no exceptional endpoint request.
\end{lemma}
\begin{proof}
Apply the classical endpoint theorem at $r+v$. The two classes of an
even-order path have equal size, so its normalized cut is $r$ and its
endpoints are in opposite classes. If their values are $E$ in the
upper class and $F$ in the lower class, summing positive differences
by vertex degrees gives
\[
 r(2r-1)=2\!\sum_{j=r+1}^{2r}j-2\!\sum_{j=1}^rj-E+F
             =2r^2-E+F.
\]
Thus $E-F=r$, which forces $F=v$.
\end{proof}

\begin{lemma}[Odd-order path interface]\label{common:odd-path}
For $u\ge2$ and $u\le X\le2u-1$, an $\alpha$-labeled path on
$\interval{1}{2u-1}$ with larger class upper, cut $u-1$, starts at
$X$ and ends at $Y=3u-1-X$, except when
\[
 2(X-1)=3(u-1).
\]
For $u=1$ the singleton word $(1)$ is allowed separately.
\end{lemma}
\begin{proof}
The two ends lie in the larger class. For $X>u$, an
$\alpha$-labeling supplied by the endpoint theorem must place that
class above the cut. If $X=u$, complementation by $2u-x$ fixes $X$
and interchanges the two classes if necessary. The forbidden upper
endpoint is exactly the displayed equality. The other endpoint is
forced by
\[
 (2u-1)(u-1)
 =2\!\sum_{j=u}^{2u-1}j-2\!\sum_{j=1}^{u-1}j-X-Y
 =2u^2-X-Y.
\]
This yields $X+Y=3u-1$.
\end{proof}

\begin{lemma}[Ordinary prescribed-first paths]\label{common:prescribed-first-path}
For every $L\ge1$ and $1\le j\le L$ there is a graceful
permutation of $\interval{1}{L}$ starting at $j$. Consequently a
single $L$-arm admits every offset $0\le k<L$.
\end{lemma}
\begin{proof}
The ordinary endpoint theorem supplies the permutation. For clarity,
even the two endpoint requests excluded only by the stronger
$\alpha$-theorem have explicit ordinary witnesses. When $L=4s+1$,
rotate the Walecki word
$(1,L,2,L-1,\ldots,2s+1)$ so that its entry at zero-based position
$2s+1$ is first. This entry is $3s+1$. The removed cut edge and the
added wrap edge both have difference $2s$, so every difference
$1,\ldots,L-1$ is retained. Complementation gives first label $s+1$.
The singleton case is immediate. To obtain an offset $k$, choose
first label $L-k$ and add $k$ to all $L$ labels. The new root edge
has color $L$, while the internal differences remain $1,\ldots,L-1$.
\end{proof}

\begin{theorem}[Three-arm hub-zero input]\label{common:hubzero-three}
\label{app:hubzero-three}
Every positive three-arm spider has an ordinary graceful labeling
with hub zero \cite[Theorem~3.2]{Panpa2025}. This input does not assert
a root-zero $\alpha$-labeling.
\end{theorem}

\subsection{Interval compositions}

\begin{lemma}[Alpha-offset amalgamation]\label{common:alpha-offset}
Let $G$ be root-zero $\alpha$ with $g$ edges and cut $K$, and let
$H$ have $h$ edges and an offset-$K$ realization. Their rooted
amalgamation is hub-zero graceful.
\end{lemma}
\begin{proof}
Keep lower labels of $G$, raise its upper labels by $h$, and keep
the offset labels of $H$. Identify roots at zero. The nonroot labels
occupy, without repetition,
\[
 \interval{1}{K},\quad\interval{K+1}{K+h},\quad
 \interval{K+h+1}{g+h}.
\]
The two palettes are $\interval{1}{h}$ and
$\interval{h+1}{h+g}$. This proves every required inventory.
\end{proof}

\begin{lemma}[Ordinary alpha amalgamation]\label{common:ordinary-amalgamation}
Let $G$ be root-zero $\alpha$ with $g$ edges and cut $K$, and let
$H$ be hub-zero graceful with $h$ edges. Their rooted amalgamation is
ordinary graceful, with final hub label $K$.
\end{lemma}
\begin{proof}
Reverse labels within the two classes of $G$ by
\[
 f^R(v)=\begin{cases}K-f(v),&f(v)\le K,\\
 g+K+1-f(v),&f(v)>K.\end{cases}
\]
Its root is $K$ and an old color $d$ becomes $g+1-d$. Raise its
uppers by $h$ and translate every label of $H$ by $K$. The common
root is $K$, and the remaining label blocks are
$\interval{0}{K-1}$, $\interval{K+1}{K+h}$ and
$\interval{K+h+1}{g+h}$. The colors are exactly
$\interval{1}{h}$ and $\interval{h+1}{g+h}$.
This is the interval form of the classical amalgamation theorem
\cite{HuangKotzigRosa1982}.
\end{proof}

\begin{lemma}[Unit stability]\label{common:unit-stability}
Arbitrarily many unit arms can be added to a hub-zero graceful
component, or to a root-zero $\alpha$ component while preserving its
index. An ordinary alpha/hub-zero decomposition therefore survives
any number of units added to the hub-zero piece before amalgamation.
\end{lemma}
\begin{proof}
For an $m$-edge component, label $v$ new leaves $m+1,\ldots,m+v$.
Their root edges have precisely these new colors. In the alpha case
each new leaf lies above the unchanged cut. The last assertion follows
by applying the preceding amalgamation to the enlarged hub-zero piece.
\end{proof}

\begin{corollary}[Ordinary hub-zero arm attachment]\label{common:ordinary-attachment}
A root-zero $\alpha$ piece of cut $K$ accepts an ordinary arm of
length $L>K$ to form a hub-zero graceful tree.
\end{corollary}
\begin{proof}
Use the single-arm offset $K$ and alpha-offset amalgamation.
\end{proof}

\begin{lemma}[Alpha-preserving arm attachment]\label{common:alpha-extension}
If a rooted tree has a root-zero $\alpha$-labeling of cut $K$, an
arm of length $L$ can be attached at the root, preserving root-zero
$\alpha$, whenever $L>2K$ and $L\ne4K+1$. The new cut is
$K+\floor{L/2}$.
\end{lemma}
\begin{proof}
Put $q=L-1$ and $\ell=\floor{(L-1)/2}$. On a separate $P_L$ choose
an $\alpha$-labeling on $\interval{0}{q}$ with lower endpoint
$j=\ell-K$. This is in the lower class because $L>2K$. The lower
exception from the endpoint theorem is exactly $L=4K+1$; the upper
exception is outside the requested lower half. Let $G$ have $g$ edges.
Keep its lower labels and raise its uppers by $L$. In the new path
map lower labels $y$ to $y+K+q-\ell+1$ and upper labels $y$ to
$y+K-\ell$. Its joining endpoint becomes $L$, and join it to the
old root zero. The four vertex blocks are
\[
 [0,K],\quad[K+1,K+q-\ell],\quad
 [K+q-\ell+1,K+q+1],\quad[K+L+1,g+L].
\]
Empty intervals are omitted. They partition $[0,g+L]$, and every
edge crosses $K+q-\ell=K+\floor{L/2}$. The new path colors are
$1,\ldots,L-1$, in reversed order; the joining color is $L$; and
the old tree supplies $L+1,\ldots,L+g$. This proves the claim.
\end{proof}

\subsection{Subdivision and the distinction between ordinary and alpha lifts}

\begin{theorem}[Uniform offset subdivision]\label{common:offset-subdivision}
If $s\in\Om(T)$ and every edge is replaced by a path of $q\ge1$
edges, then
\[
 \interval{qs}{q(s+1)-1}\ \subseteq\ \Om(T[q]).
\]
The complete column construction and signed-difference proof are given
in Section~\ref{app:subdivision}. The input need not be alpha.
\end{theorem}

In particular ordinary gracefulness is preserved under uniform
subdivision: root an ordinary labeling at its zero-labeled vertex and
take $s=0$. That vertex need not be the spider's hub. The underlying
graph operation multiplies all arm lengths by $q$, independently of
which vertex was used as the root during construction.

For alpha preservation a stronger hypothesis is used. Section~
\ref{app:lifts} proves balanced-index alpha subdivision and the ordinary
perfect-matching lift, with their exact root and index conditions.
An arbitrary alpha labeling is never silently assumed to survive
subdivision as an alpha labeling.
